\chapter{Related Work}
\label{chapter:related_work}

This chapter reviews four bodies of literature: AI agent architectures (\cref{sec:rw_agents}), the attack and defence landscape for LLM-based systems (\cref{sec:rw_security}), existing benchmarks for agent security and privacy (\cref{sec:rw_benchmarks}), and trust and access control models (\cref{sec:rw_trust}). We conclude by positioning SharedOS at the intersection (\cref{sec:rw_position}).


% ============================================================================
\section{AI Agent Architectures}
\label{sec:rw_agents}
% ============================================================================

\textbf{Single-agent systems.} ReAct~\cite{yao2023react} established the reasoning-action loop that underpins modern agents. Toolformer~\cite{schick2023toolformer} showed tool use can be a learned competency, while WebShop~\cite{yao2022webshop}, WebArena~\cite{zhou2024webarena}, SWE-bench~\cite{jimenez2024swebench}, and AgentBench~\cite{liu2024agentbench} evaluate agents in interactive environments that require tool use, long-horizon decision making, and task completion rather than isolated text generation. MemGPT~\cite{packer2024memgpt} introduced an OS-inspired memory hierarchy (main context, archival storage, paging), transforming stateless LLMs into persistent agents, and Generative Agents~\cite{park2023generative} showed how persistent memory and reflection can support believable social behaviour. These advances have translated into commercial systems such as Devin~\cite{devin2024} (autonomous software engineering) and Manus~\cite{manus2025} (desktop automation with GUI control), alongside research systems such as ChemCrow~\cite{bran2024chemcrow} and R\&D-Agent-Quant~\cite{li2025rdagentquantmultiagentframeworkdatacentric}, which apply tool-using agents to scientific and quantitative R\&D workflows. Individual agents are now capable enough to serve as personal delegates; the bottleneck is secure cross-boundary coordination.

\textbf{Multi-agent frameworks.} AutoGen~\cite{wu2023autogen} enables multi-agent conversation through structured role-based patterns. CrewAI~\cite{crewai2024} provides role-based workflows with defined goals and tool access. MetaGPT~\cite{hong2023metagpt} and ChatDev~\cite{qian2023chatdev} simulate software organisations with specialised roles, encoding standard operating procedures as structured communication protocols. CAMEL~\cite{li2023camel} explores emergent communication between role-playing agents. The critical observation is that every existing multi-agent system is \textit{multi-role-within-one-owner}: all agents share a single principal. None address the scenario where agents representing different owners with potentially conflicting privacy interests must coordinate.

\textbf{Agent operating systems.} AIOS~\cite{mei2025aios} proposes OS-level primitives (scheduler, context manager, tool manager) for managing multiple agents on shared infrastructure, but its isolation is resource-oriented (preventing starvation), not privacy-oriented (preventing information leakage). OS-Copilot~\cite{wu2024oscopilot} operates \textit{within} the OS as a single-owner agent with no cross-boundary considerations. No existing agent OS addresses information flow control across trust boundaries. SharedOS occupies this gap.


% ============================================================================
\section{LLM Security and Privacy}
\label{sec:rw_security}
% ============================================================================

\textbf{Attacks.} Zou et al.~\cite{zou2023universal} showed gradient-based adversarial suffixes (GCG) break RLHF alignment and transfer across models. AutoDAN~\cite{liu2024autodan} generates fluent jailbreaks via genetic algorithms, and PAIR~\cite{chao2025pair} uses an attacker LLM to refine prompts in under 20 queries. PAP~\cite{zeng2024pap} applies 40 persuasion techniques for $>$90\% attack success on GPT-4, directly relevant to agent-to-agent settings. Indirect prompt injection further shows that tool-using applications blur the boundary between data and instructions: retrieved content can become an attack vector even when the user prompt itself is benign~\cite{greshake2023not}. Crescendo~\cite{russinovich2025crescendo} demonstrates multi-turn escalation through individually innocuous turns. Most critically, Nasr et al.~\cite{nasr2025attacker} tested twelve published defences under adaptive attacks and found that \textit{every defence} could be bypassed with $>$90\% success. This motivates our architectural approach: security must be enforced by controlling what information is \textit{present} rather than instructing the model not to reveal it.

\textbf{Defences.} Wallace et al.~\cite{wallace2024instruction} proposed hierarchical instruction prioritisation (system $>$ user $>$ tool). Hines et al.~\cite{hines2024spotlighting} introduced data-marking transformations that help models distinguish instructions from data, and StruQ~\cite{chen2025struq} formalises a related separation by placing instructions and data into distinct structured channels. Llama Guard~\cite{inan2023llamaguard} and NeMo Guardrails~\cite{rebedea2023nemo} provide classifier-based filtering. All operate at the prompt or model-interface level: necessary but insufficient under adaptive attack~\cite{nasr2025attacker}. This motivates both our defence gradient methodology and the Mountable Context Cell design.

\textbf{Privacy in Language Models} Dwork et al.~\cite{dwork2006differential} established differential privacy (DP); Yang et al.~\cite{yang2026dpagent} applied DP to agent communication, adding calibrated noise to bound per-query leakage at a utility cost our work measures empirically. Carlini et al.~\cite{carlini2021extracting} demonstrated that LLMs memorise and reproduce verbatim training data, establishing that LLMs are fundamentally unreliable as privacy-preserving components. Any system depending on LLM behaviour for privacy inherits this unreliability.

% ============================================================================
\section{Agent Security Benchmarks}
\label{sec:rw_benchmarks}
% ============================================================================

A growing number of benchmarks evaluate agent security. We review them along six dimensions: cross-boundary interaction, tool-mediated evaluation, dual-metric measurement, multi-turn scenarios, relationship modelling, and action safety.

\textbf{Prompt injection benchmarks.} InjecAgent~\cite{zhan2024injecagent} evaluates indirect injection in tool-integrated agents (1,054 tests, 17 tools) but within single-agent contexts without measuring utility cost. AgentDojo~\cite{debenedetti2024agentdojo} provides dual-metric evaluation (97 tasks, 629 security tests) but operates within a single task context with no cross-boundary interaction. TensorTrust~\cite{toyer2024tensortrust} gamified prompt injection (126,000 attack/defence pairs), demonstrating human creativity but not modelling agent-to-agent settings. HackAPrompt~\cite{schulhoff2024hackaprompt} confirmed through 600,000+ adversarial prompts that human creativity consistently defeats automated defences.

\textbf{Multi-agent and cross-boundary benchmarks.} AgentSocialBench~\cite{wang2026agentsocialbench} evaluates privacy in agentic social networks, discovering the ``abstraction paradox'' (more capable agents are more susceptible to social engineering) but does not evaluate defences or measure utility. ConVerse~\cite{gomaa2025converse} is the most comparable setup: personal assistants interacting with external providers (864 attack scenarios, up to 88\% success), but measures attack success only with no utility dimension. PAC-BENCH~\cite{park2026pacbench} evaluates collaboration under privacy constraints with dual measurement but not across genuine trust boundaries. MAGPIE~\cite{juneja2025magpie} provides 200 collaborative tasks but all agents share a single principal. AgentLeak~\cite{yagoubi2026agentleak} identifies seven distinct leakage channels but is diagnostic only. Agents of Chaos~\cite{shapira2026chaos} red-teams deployed systems, demonstrating lab-to-production attack transfer without systematic evaluation methodology.

\Cref{table:benchmark_comparison} summarises these benchmarks across the six dimensions critical to cross-boundary privacy evaluation.

\begin{table}[ht]
\centering
\caption{Comparison of agent security benchmarks. $\checkmark$\ = fully supported, $\sim$ = partially supported, -- = not supported.}
\label{table:benchmark_comparison}
\resizebox{\textwidth}{!}{%
\begin{tabular}{lcccccc}
\toprule
\textbf{Benchmark} & \textbf{Cross-boundary} & \textbf{Tool-mediated} & \textbf{Dual metric} & \textbf{Multi-turn} & \textbf{Relationship} & \textbf{Action safety} \\
\midrule
InjecAgent \cite{zhan2024injecagent}          & --           & $\checkmark$ & --           & --           & --           & $\checkmark$ \\
AgentDojo \cite{debenedetti2024agentdojo}     & --           & $\checkmark$ & $\sim$       & --           & --           & $\sim$       \\
TensorTrust \cite{toyer2024tensortrust}       & --           & --           & --           & --           & --           & --           \\
HackAPrompt \cite{schulhoff2024hackaprompt}   & --           & --           & --           & --           & --           & --           \\
AgentSocialBench \cite{wang2026agentsocialbench} & $\checkmark$ & --       & --           & --           & --           & --           \\
ConVerse \cite{gomaa2025converse}             & $\checkmark$ & --           & --           & $\checkmark$ & --           & --           \\
PAC-BENCH \cite{park2026pacbench}            & $\sim$       & $\checkmark$ & $\checkmark$ & --           & --           & --           \\
MAGPIE \cite{juneja2025magpie}               & --           & $\checkmark$ & --           & --           & --           & $\sim$       \\
AgentLeak \cite{yagoubi2026agentleak}        & $\checkmark$ & $\checkmark$ & --           & --           & --           & --           \\
Agents of Chaos \cite{shapira2026chaos}      & $\checkmark$ & $\checkmark$ & --           & $\sim$       & --           & $\checkmark$ \\
\midrule
\textbf{PACT (Ours)} & $\checkmark$ & $\checkmark$ & $\checkmark$ & $\checkmark$ & $\checkmark$ & $\checkmark$ \\
\bottomrule
\end{tabular}%
}
\end{table}

\textbf{No existing benchmark combines all six dimensions.} InjecAgent and AgentDojo measure tool-mediated attacks within single-agent contexts. ConVerse crosses trust boundaries but measures only attack success. PAC-BENCH provides dual measurement but not across genuine trust boundaries. None model how the relationship between requester and data owner affects the security-utility frontier. PACT fills this gap.


% ============================================================================
\section{Trust and Access Control}
\label{sec:rw_trust}
% ============================================================================

Classical access control operates on a principle where policy \textit{is} enforcement. RBAC assigns permissions to predefined roles; ABAC evaluates access based on subject/object/environment attributes; Bell-LaPadula~\cite{bell1973secure} formalises mandatory access control (``no read up, no write down'') for confidentiality; and Denning's lattice model~\cite{denning1976lattice} constrains information flow through a partial order on security classes. All assume deterministic enforcement and static classification hierarchies.

Capability-based security~\cite{dennis1966programming} takes a different approach: unforgeable tokens grant specific permissions, enforcing the principle of least authority. A process receives only the capabilities it needs. Our Mountable Context Cell design is directly inspired by this model: an agent receives only the context cells relevant to the current interaction and cannot access data outside those cells regardless of its reasoning.

The key distinction in our setting is that access control is \textit{relationship-relative}: the same information may be shareable with one requester but not another. Moreover, the ``policy'' (a system prompt or permission instruction) is not enforcement. It is \textit{input to reasoning}. Under adversarial conditions, the model can reason its way around any natural-language instruction. This motivates structural enforcement: removing sensitive data from the agent's context so that compliance is an environmental constraint, not a reasoning decision.


% ============================================================================
\section{Position: SharedOS at the Intersection}
\label{sec:rw_position}
% ============================================================================

The realistic interaction of agents as personal delegates is not well modelled by any existing work, because each delegate holds private state on behalf of a distinct owner. Multi-agent frameworks assume shared principals. Security benchmarks measure attack success without a utility dimension. No existing benchmark combines cross-boundary interaction, tool mediation, dual metrics, multi-turn scenarios, relationship modelling, and action safety (\Cref{table:benchmark_comparison}).

This thesis formalises the cross-boundary delegation problem, proposes SharedOS as a platform that instantiates it with persistent state, tool access, and social relationships, and investigates the security--utility frontier through systematic experimentation (PACT-PAIR, PACT-NET). Where prompt-level defences reach their ceiling, we propose two architectural solutions, Mountable Context Cells and the Escalation Protocol, that move enforcement from reasoning to structure.
